\section{Conclusion}
We introduce DiSPaT, a self-play fine-tuning framework for unsupervised tabular anomaly detection with LLMs. Our method uses an iterative refinement mechanism grounded in $f$-divergence minimization: at each iteration, the model generates synthetic samples that serve as pseudo-anomalies, and a discriminative objective drives the policy to better align with the true normal data distribution. This self-play loop provides the discriminative signal absent in one-shot fine-tuning approaches, enabling the model to learn a tighter description of normality. Theoretically, we establish a tight variational approximation connecting our training objective to the $f$-divergence between real and model-induced distributions, providing a principled foundation for iterative self-improvement. Empirically, DiSPaT consistently outperforms AnoLLM, strong classical and deep-learning baselines, and additional recent tabular anomaly detection methods across multiple datasets, supporting the effectiveness of the proposed framework for tabular anomaly detection.